% Diagrama 5. Accionamiento del actuador lineal (etapas 2 y 3), en TikZ.
% Equivale a D5_accionamiento_actuador.mmd. (P) = propuesto.
% Compilar: pdflatex D5_accionamiento_actuador_tikz.tex
\documentclass[border=4pt]{standalone}
\usepackage[T1]{fontenc}
\usepackage[utf8]{inputenc}
\usepackage[spanish,es-noshorthands]{babel}
\usepackage{times}
\usepackage{amssymb}
\usepackage{tikz}
% Estilos comunes de los diagramas de validación (TikZ).
% Blanco y negro, Times a 8 pt, sin colores, como la Figura 22 de la tesis.
\usetikzlibrary{positioning,calc,shapes.geometric}
\tikzset{
  base/.style   = {font=\fontsize{8}{9.6}\selectfont},
  proc/.style   = {draw, fill=white, align=center, text width=3.0cm,
                   minimum height=0.9cm, inner sep=3pt,
                   execute at begin node={\hyphenpenalty=10000\relax}},
  propp/.style  = {proc, dashed},
  negra/.style  = {proc, fill=black!12},
  falla/.style  = {proc, fill=black!7},
  dec/.style    = {draw, fill=white, diamond, aspect=1.9, align=center,
                   inner sep=0.5pt, text width=1.95cm,
                   execute at begin node={\hyphenpenalty=10000\relax}},
  term/.style   = {draw, fill=black!4, rounded corners=0.42cm, align=center,
                   text width=2.7cm, minimum height=0.84cm, inner sep=3pt,
                   execute at begin node={\hyphenpenalty=10000\relax}},
  con/.style    = {draw, circle, fill=white, minimum size=0.6cm, inner sep=0pt},
  nota/.style   = {draw, dashed, fill=white, align=center, text width=4.6cm,
                   inner sep=3pt, font=\fontsize{7.2}{8.6}\selectfont},
  marco/.style  = {draw, rounded corners=5pt, black!55, thin},
  tmarco/.style = {font=\fontsize{8}{9.6}\selectfont\itshape, anchor=north west},
  fl/.style     = {->, >=latex, semithick},
  ln/.style     = {semithick},
  et/.style     = {font=\fontsize{6.8}{8}\selectfont, inner sep=1.2pt, fill=white},
  union/.style  = {circle, fill=black, inner sep=0pt, minimum size=3pt}
}


\begin{document}
\begin{tikzpicture}[base, bl/.style={proc, minimum height=1.3cm}]
% camino de mando
\node[bl, text width=2.35cm] (PI)  at (0,0)     {\textbf{Raspberry Pi Zero 2 W}\\orden de flexión con confianza alta o autorización};
\node[bl, text width=2.6cm]  (MCU) at (3.95,0)  {\textbf{ESP32 DevKit}\\extensión hasta 35 mm por posición (6.9 s); corte si no avanza o a los 12 s};
\node[bl, text width=2.2cm]  (DRV) at (7.85,0)  {\textbf{DRV8871}\\módulo, puente completo};
\node[bl, text width=2.9cm]  (ACT) at (11.75,0) {\textbf{Actuonix L12-50-210-12-P}\\62 N en máxima potencia, 5.1 mm/s en servicio, carrera 50 mm};
\node[bl, text width=2.1cm]  (LAM) at (15.45,0) {\textbf{Lámina PEI}\\flecha de trabajo 30 mm};

\draw[fl] (PI) -- node[et, above] {UART} node[et, below] {TX2/RX2} (MCU);
\draw[fl] (MCU) -- node[et, above] {D25: IN1} node[et, below] {D26: IN2} (DRV);
\draw[fl] (DRV) -- node[et, above] {OUT1, OUT2} node[et, below] {por J4} (ACT);
\draw[fl] (ACT) -- node[et, above] {vástago} (LAM);
\draw[fl, dashed] ($(ACT.north)+(0.6,0)$) -- ++(0,0.45) -| node[et, above, pos=0.1] {potenciómetro (3.3 V): cursor a D34} ($(MCU.north)+(0.6,0)$);

% limitación de corriente, pendiente
\node[negra, text width=4.2cm] (R3) at (7.85,2.55) {R3 56 k$\Omega$ en ILIM: solo respaldo ante un cortocircuito; el bloqueo se detecta por la posición};
\draw[fl, dashed] (R3) -- (DRV);

% alimentación y botón de emergencia
\node[bl, text width=2.35cm] (SETA) at (0,-2.6)  {\textbf{Botón de emergencia NC} (J2)\\ISO 13850, apertura positiva};
\node[bl, text width=2.0cm]  (BUS)  at (3.95,-2.6) {\textbf{+24V\_ACT}\\riel tras el botón de emergencia};
\node[bl, text width=2.2cm]  (BUCK) at (7.85,-2.6) {\textbf{LM2596}\\reductor a 12 V};
\draw[fl] (SETA) -- (BUS);
\draw[fl] (BUS) -- (BUCK);
\draw[fl] (BUCK) -- node[et, right=1pt, align=left] {VM\\C3 100 $\mu$F} (DRV);
\draw[fl, dashed] (SETA.north) -- node[et, left=1pt, pos=0.5, align=right] {divisor R7/R8\\a D35} ++(0,0.9) -| ($(MCU.south)+(-0.6,0)$);

% tabla de verdad
\node[nota, text width=5.0cm] at (13.3,-2.6) {\textbf{Tabla de verdad del DRV8871}\\IN1 = 1, IN2 = 0: un sentido\\IN1 = 0, IN2 = 1: sentido opuesto\\0, 0: libre;\ \ 1, 1: freno\\Sentido de extensión por fijar en el montaje};
\end{tikzpicture}
\end{document}
